\chapter{اللقاء الأول: مقدمة في دراسة تفسير الطبري}

السلام عليكم ورحمة الله وبركاته.

الحمد لله رب العالمين، وأشهد أن لا إله إلا الله وحده لا شريك له، وأشهد أن محمداً عبده ورسوله. اللهم صل على محمد وعلى آل بيته وذريته كما صليت على آل إبراهيم إنك حميد مجيد، اللهم بارك على محمد وعلى آل بيته وذريته كما باركت على آل إبراهيم إنك حميد مجيد.

\begin{ayah}
يَا أَيُّهَا الَّذِينَ آمَنُوا اتَّقُوا اللَّهَ وَلْتَنظُرْ نَفْسٌ مَّا قَدَّمَتْ لِغَدٍ ۖ وَاتَّقُوا اللَّهَ ۚ إِنَّ اللَّهَ خَبِيرٌ بِمَا تَعْمَلُونَ * وَلَا تَكُونُوا كَالَّذِينَ نَسُوا اللَّهَ فَأَنسَاهُمْ أَنفُسَهُمْ ۚ أُولَٰئِكَ هُمُ الْفَاسِقُونَ * لَا يَسْتَوِي أَصْحَابُ النَّارِ وَأَصْحَابُ الْجَنَّةِ ۚ أَصْحَابُ الْجَنَّةِ هُمُ الْفَائِزُونَ * لَوْ أَنزَلْنَا هَٰذَا الْقُرْآنَ عَلَىٰ جَبَلٍ لَّرَأَيْتَهُ خَاشِعًا مُّتَصَدِّعًا مِّنْ خَشْيَةِ اللَّهِ ۚ وَتِلْكَ الْأَمْثَالُ نَضْرِبُهَا لِلنَّاسِ لَعَلَّهُمْ يَتَفَكَّرُونَ
\end{ayah}

قال الله تبارك وتعالى:
\begin{ayah}
يَا أَيُّهَا النَّاسُ قَدْ جَاءَتْكُم مَّوْعِظَةٌ مِّن رَّبِّكُمْ وَشِفَاءٌ لِّمَا فِي الصُّدُورِ وَهُدًى وَرَحْمَةٌ لِّلْمُؤْمِنِينَ * قُلْ بِفَضْلِ اللَّهِ وَبِرَحْمَتِهِ فَبِذَٰلِكَ فَلْيَفْرَحُوا هُوَ خَيْرٌ مِّمَّا يَجْمَعُونَ
\end{ayah}

وفضل الله تبارك وتعالى في قول كثير من المفسرين هو الإسلام، ورحمته القرآن الذي جعله الله تبارك وتعالى نوراً نحيا به ونهتدي ونمشي به في الناس. والحمد لله رب العالمين على نعمة الإسلام، والحمد لله الذي أنزل على عبده الكتاب. 

وأحمد الله تبارك وتعالى أن جمعني مع الإخوة الكرام في مدارسة العلوم الشرعية والفقه في الدين والاجتماع على ذكر الله تبارك وتعالى. وفي مسند أحمد وصحيح مسلم:

\begin{hadith}
\isnad{من حديث} \rawi{أبي إسحاق} \isnad{عن} \rawi{الأغر أبي مسلم}، \isnad{أشهد على} \rawi{أبي هريرة} \rawi{وأبي سعيد} \isnad{أنهما شهدا على النبي} صلى الله عليه وسلم أنه قال:

لا يقعد قوم يذكرون الله عز وجل إلا حفتهم الملائكة، وغشيتهم الرحمة، ونزلت عليهم السكينة، وذكرهم الله فيمن عنده.
\end{hadith}

الاجتماع على الفقه في الدين من خير ما يذكر به الله تبارك وتعالى. اللهم أنزل علينا السكينة وارحمنا واذكرنا فيمن عندك، وأصلي على عبدالله ورسوله محمد بن عبدالله الذي مَنّ الله تبارك وتعالى على المؤمنين ببعثته. قال الله عز وجل:

\begin{ayah}
لَقَدْ مَنَّ اللَّهُ عَلَى الْمُؤْمِنِينَ إِذْ بَعَثَ فِيهِمْ رَسُولًا مِّنْ أَنفُسِهِمْ يَتْلُو عَلَيْهِمْ آيَاتِهِ وَيُزَكِّيهِمْ وَيُعَلِّمُهُمُ الْكِتَابَ وَالْحِكْمَةَ وَإِن كَانُوا مِن قَبْلُ لَفِي ضَلَالٍ مُّبِينٍ
\end{ayah}

قال الله تبارك وتعالى:
\begin{ayah}
إِنَّ اللَّهَ وَمَلَائِكَتَهُ يُصَلُّونَ عَلَى النَّبِيِّ ۚ يَا أَيُّهَا الَّذِينَ آمَنُوا صَلُّوا عَلَيْهِ وَسَلِّمُوا تَسْلِيمًا
\end{ayah}

اللهم صل على محمد وعلى آل محمد كما صليت على آل إبراهيم إنك حميد مجيد، اللهم بارك على محمد وعلى آل محمد كما باركت على آل إبراهيم إنك حميد مجيد.

اللهم ارض عن صحابة نبيك الذين آمنوا به وعزروه ونصروه واتبعوا النور الذي أنزل معه، وكانوا خير أمة أخرجت للناس، يأمرون بالمعروف وينهون عن المنكر ويسارعون في الخيرات، حفظوا لنا هدي رسول الله صلى الله عليه وسلم وسننه وبلغوها وبينوها.

ورحم الله أهل العلم في كل زمان ومكان، الذين يمسكون بالكتاب ويدعون إلى الله ويبينون الهدى، احتسبوا أعمارهم لطلب العلم ونصر الله ورسوله وكتابه، فاللهم اجزهم عنا خيراً. وإني أحمد الله تبارك وتعالى حمداً كثيراً طيباً مباركاً فيه أن جمعني وإياكم على التماس علوم الشريعة والتفقه في الدين، والاشتغال بطلب القرآن وبيانه من حديث رسول الله صلى الله عليه وسلم وما كان عليه السابقون الأولون الذين هم خير أمة أخرجت للناس.

ومما ينبغي أن تعلمه وأن يكون ثابتاً في قلبك أن أصل الهدى، وأن دين الحق، وأن العلم النافع والعمل الصالح، هو الإيمان بالله ورسله. والله تبارك وتعالى قال:
\begin{ayah}
يَا بَنِي آدَمَ إِمَّا يَأْتِيَنَّكُمْ رُسُلٌ مِّنكُمْ يَقُصُّونَ عَلَيْكُمْ آيَاتِي ۙ فَمَنِ اتَّقَىٰ وَأَصْلَحَ فَلَا خَوْفٌ عَلَيْهِمْ وَلَا هُمْ يَحْزَنُونَ * وَالَّذِينَ كَذَّبُوا بِآيَاتِنَا وَاسْتَكْبَرُوا عَنْهَا أُولَٰئِكَ أَصْحَابُ النَّارِ ۖ هُمْ فِيهَا خَالِدُونَ
\end{ayah}

ومن آمن بالله ورسله واتبع الهدى فهو من خير البرية. ومحمد بن عبدالله بن عبد المطلب هو خاتم النبيين، ودينه الإسلام هو الدين الحق، ولن يهتدي أحد إلا بقدر علمه ويقينه وصبره على هذا الدين، وبقدر علمه بالكتاب الذي أنزله الله تبارك وتعالى للناس كافة بشيراً ونذيراً، وبقدر الاهتداء بهدي النبي صلى الله عليه وسلم. قال الله تبارك وتعالى:
\begin{ayah}
وَالَّذِينَ آمَنُوا وَعَمِلُوا الصَّالِحَاتِ وَآمَنُوا بِمَا نُزِّلَ عَلَىٰ مُحَمَّدٍ وَهُوَ الْحَقُّ مِن رَّبِّهِمْ ۙ كَفَّرَ عَنْهُمْ سَيِّئَاتِهِمْ وَأَصْلَحَ بَالَهُمْ
\end{ayah}

وكان النبي صلى الله عليه وسلم يفتتح خطبه فيقول:
\begin{hadith}
فإن خير الحديث كتاب الله، وخير الهدي هدي محمد، وشر الأمور محدثاتها، وكل بدعة ضلالة.
\end{hadith}

فأصدق العلم وأحسن الكلام هو القرآن الكريم. وأحسن الهدي - يعني العمل والأخلاق - هدي رسول الله صلى الله عليه وسلم. وكل إنسان في هذه الدنيا له قصد وعمل، له غاية يطلبها ووسيلة إلى هذه الغاية، فمن قصد هدى رسول الله صلى الله عليه وسلم وعمل به اهتدى، ومن لم يقصده أو لم يعمل به ضل وغوى.

القرآن الكريم كما سبق كثيراً في الحديث عنه، هو نور يتصور به المسلم الحياة ويحسن التصرف فيها بقدر علمه وعمله. قال الله تبارك وتعالى:
\begin{ayah}
أَوَمَن كَانَ مَيْتًا فَأَحْيَيْنَاهُ وَجَعَلْنَا لَهُ نُورًا يَمْشِي بِهِ فِي النَّاسِ كَمَن مَّثَلُهُ فِي الظُّلُمَاتِ لَيْسَ بِخَارِجٍ مِّنْهَا ۚ كَذَٰلِكَ زُيِّنَ لِلْكَافِرِينَ مَا كَانُوا يَعْمَلُونَ
\end{ayah}

فكتاب الله تبارك وتعالى وبيانه من حديث رسول الله صلى الله عليه وسلم هما خير ما يمكن أن يشتغل به إنسان في هذه الحياة الدنيا، \begin{ayah}رَّسُولٌ مِّنَ اللَّهِ يَتْلُو صُحُفًا مُّطَهَّرَةً\end{ayah}.

أعظم ما ينبغي أن تطلبه نفس: العلم بربها، والعلم بوحي الله وما خلقت له. وأعظم ما تشتغل به في هذه الحياة الدنيا الاستقامة على دين الله تبارك وتعالى. وأعقل الناس وأذكاهم وأزكاهم وأحسنهم عملاً هم الذين يطلبون كتاب الله وهدي رسول الله صلى الله عليه وسلم، والخاسرون أعمالاً الذين عاشوا في هذه الدنيا لم يطلبوا وحي الله ولم يتعلموا هدي رسول الله صلى الله عليه وسلم.

وبعد، ففي هذا اليوم، يوم الأربعاء بعد فجر الأول من شهر صفر من عام ثلاث وأربعين وأربعمائة وألف من هجرة رسول الله صلى الله عليه وسلم، ومحدثكم أخوكم حسين بن عبدالرازق المصري، وهذا هو لقاؤنا الأول في الاجتماع على دراسة تفسير ابن جرير الطبري رحمه الله.

وتأتي هذه الدورة حلقة جديدة من دورات دراسة تراث الأئمة المحققين، والتي بدأتها مع طلاب العلم قديماً قبل عشرين عاماً. ولكني بعد انتشار وسائل التواصل ووجود عدد كثير من طلاب العلم يعتنون بهذه المواقع، أحببت أن يكون لنا بعض الدورات النوعية في هذا الباب.

وتذكرت هذا الأثر لحبان بن موسى، وهو أحد تلامذة عبد الله بن المبارك رحمه الله، قال: عوتب عبدالله بن المبارك فيما يفرق من المال في البلدان (يعني كان ينفق على كثير من طلاب العلم في البلدان المختلفة فعوتب على ذلك)، فقال: \textit{«إني لأعرف مكان قوم لهم فضل وصدق طلبوا الحديث فأحسنوا طلبه، والناس محتاجون إليهم وهم بحاجة إلى أنفسهم وذراريهم، فإن تركناهم ضاع علمهم، وإن أعناهم بثوا العلم لأمة محمد صلى الله عليه وسلم، لا أعلم بعد النبوة أفضل من بث العلم»}.

فأنا أقول: أنا أعلم كثيراً من طلبة العلم في مختلف بلدان العالم، يحبون الله ورسوله فيما أحسب، ويطلبون كتاب الله وحديث رسول الله صلى الله عليه وسلم، يريدون التفقه في الدين. وأكثر هؤلاء الشباب عنده حرص وجلد وصبر، يجاهد نفسه ويقاوم شهواته ويلقي بنفسه في معالي الأمور، يبذل من وقته وماله وتفكيره، يريد نصر الإسلام، يريد أن يكون له دور في سفينة الإسلام، وأن يكون حلقة في سلسلة التعلم والتعليم. تتنوع أعمال هؤلاء وسبلهم في نصرة الإسلام والغاية واحدة. وأنا أعرف كثيراً جداً من هؤلاء الشباب، وأكثر من الدعاء لهم والتشجيع وأفخر بهم، وأرجو أن يبلغهم الله تبارك وتعالى من الخير ما يطلبون.

وأرجو من الله أن أكون عوناً لهم ومشجعاً وموجهاً على قدر ما عندي من العلم. وإني والله لأضن بتلك الجهود والأوقات أن تضيع في غير الطريق الصحيح، أو أن تتعطل المواهب وتختفي انشغالاً بما لا يستحق. طالب العلم كما يحتاج إخلاصاً وعزماً وصبراً وجلداً، فإنه كذلك يحتاج توجيهاً وإرشاداً حتى لا تهدر قدرته أو أوقاته في الطريق الخطأ. فأبدأ مستعيناً بالله تبارك وتعالى، سائلاً توفيقه وهداه لي وللمشاركين الكرام: اللهم اهدنا وسددنا، اللهم هيئ لنا من أمرنا رشداً واربط على قلوبنا، اللهم اجعلنا من أهل الصدق والبيان، اللهم لا تجعلنا بما تعلمنا ظهيراً للمجرمين.

\section*{تفاصيل الدورة ومقاصدها}
دورات دراسة الكتب التأسيسية في علوم الشريعة، هذه حلقة منها: دورة دراسة \textbf{«جامع البيان عن تأويل آي القرآن»} وهو المشهور بتفسير الطبري. والطبري هو أبو جعفر محمد بن جرير الطبري المتوفى عام 310 هـ.

\subsection*{المقاصد العامة للدورة}
\begin{enumerate}
    \item التشجيع على قراءة المطولات من كتب أهل العلم والمدارسة الجماعية.
    \item التدريب على دراسة كتب الأئمة المحققين (القراءة والضبط والفهم).
    \item التعرف على أبواب علوم القرآن وتفسيره مما يثيره الطبري أثناء التفسير.
    \item التدريب على استخراج الفوائد وتصنيفها تحت بابها اللائق.
    \item الوقوف على الأمثلة التطبيقية تحت كل باب من أبواب علوم القرآن والتفسير، فبدراستنا لتفسير الطبري ستقف على مئات الأمثلة إن شاء الله.
    \item جمع النظائر والشواهد والأمثلة تحت الباب لتكون مادة لمن يريد أن يدرس هذا العلم.
    \item أن يكون الطالب جزءاً من الدرس ومشاركاً فيه، بالتحضير قبله، وبالمناقشة أثنائه، وكذلك بعد الدرس بالمراجعة واستخراج الفوائد وتصنيفها.
\end{enumerate}

\subsection*{المستهدف من هذه الدورات}
\begin{itemize}
    \item \textbf{طلبة العلم عموماً:} لأن دراستنا لهذا الكتاب ليس المراد منها دراسة تفسير القرآن كعلم فحسب؛ بل نحن نريد أن يكون تفسير القرآن ضمن الأهداف المطلوبة، لكن المقصد الأساس هو التدريب على قراءة تراث الأئمة المحققين. نحن ندرس تراث هؤلاء لننتفع منه من كل جانب، نستنطق هذا الكتاب لنستخرج ما فيه من أبواب العلم والهدى والنظر والنقد والمناقشة والتقرير.
    \item \textbf{طلبة علوم القرآن بشكل خاص:} الذين يعتنون بكل علم يتصل بكتاب الله، فكتاب الطبري هو كنز تطبيقي في علوم القرآن وبشكل أخص في علم التفسير.
\end{itemize}

والذي يريد أن يدرس هذا الصنف من الكتب لابد أن يكون عنده أمران:
\begin{enumerate}
    \item \textbf{التأهب للاستفادة:} لأن هذه الكتب تدخل في كثير من أبواب العلم، فيها الكلام عن علوم القرآن والحديث والأسانيد والفقه ومسائل الإيمان والفرق واللغة والشعر، فأنت ستدخل على بحر من العلوم يحتاج إلى استعداد.
    \item \textbf{خطة للدراسة:} لأن هذا الكتاب كبير (حوالي 26 مجلداً)، وهذه الكتب المطولات كثيراً ما تُقرأ مرة واحدة في العمر، فتحتاج إلى خطة لتكون كشافاً لك في دراسة هذا الكتاب وتتعلم كيف تستخرج أقصى ما يمكن من علومه.
\end{enumerate}

\subsection*{خطة دراستنا للكتاب}
الدروس ستكون مباشرة إن شاء الله تبارك وتعالى عبر قناة التليجرام يومي الأربعاء والخميس من كل أسبوع. هذا هو الأصل، إلا أن يحصل ظرف فأغير الموعد. مدة الدرس ستكون من ثلاث إلى أربع ساعات لأن الكتاب كبير، ولا نستطيع أن نقلل عن ثلاث ساعات، يعني في الأسبوع ست ساعات.

\section*{نشأة علوم الشريعة وتصورها}
أحب أن أبين قبل الدخول في دراسة تفسير الطبري، أن تتصور علوم الشريعة؛ لأن كثيراً من الطلاب لا يعرف كيف بدأت هذه العلوم. لقد استقر الآن عنده وجود علم التفسير وعلوم القرآن، علوم الحديث، مصطلح الحديث، علم الرجال، العلل، تخريج الحديث، علم الفقه، أصول الفقه، السيرة النبوية، والتاريخ... ولكنه لا يعرف كيف بدأت ومم نشأت. أنا أذكر لك هنا مختصراً لتتصور لماذا نحن ندرس هذه الكتب؟ ولماذا نعتني بكتب هؤلاء العلماء؟

خلاصة علوم الشريعة أنها لا تخرج عن أمرين: إما \textbf{نقل ورواية}، وإما \textbf{فقه وفهم وتفسير}. كل العلوم في الدنيا عموماً إما نقل لما سبق، أو فهم لهذا المنقول واستخراج للفوائد منه والبناء عليه.

\begin{itemize}
    \item \textbf{الصنف الأول (النقل والرواية):} وهي المعطيات الموجودة: القرآن الكريم، وحديث رسول الله صلى الله عليه وسلم، وآثار الصحابة، ثم آثار التابعين. هذه المرويات هي النقل. وفي هذا الباب يحتاج الطالب أن يتعلم مصادر العلم التي حفظت لنا هذه النصوص، ويحتاج أن يتعلم كيف يجمع وكيف ينقد هذه الآثار ليستخرج أصحها وأحسنها.
    \item \textbf{الصنف الثاني (الفقه والفهم والتفسير):} وهو جودة التفاعل مع هذه المعطيات المنقولة.
\end{itemize}

وهذا التقسيم هو ما ذكره النبي صلى الله عليه وسلم في حديث من جوامع الكلم، حيث قال:
\begin{hadith}
لِيُبَلِّغِ الشَّاهِدُ مِنْكُمُ الغَائِبَ، فَإِنَّ الشَّاهِدَ عَسَى أَنْ يُبَلِّغَ مَنْ هُوَ أَوْعَى لَهُ مِنْهُ.
\end{hadith}

وكذلك في حديث \rawi{أبي جحيفة وهب بن عبدالله السوائي}، لما سأل \rawi{علي بن أبي طالب} رضي الله عنه: «هل عندكم شيء من الوحي إلا ما في كتاب الله؟»، قال:
\begin{hadith}
لا، والذي فلق الحبة وبرأ النسمة، ما أعلمه إلا فهماً يعطيه الله رجلاً في القرآن.
\end{hadith}

هذه الأحاديث تبين خلاصة ما يحتاجه طالب الفقه في الدين: العناية بالروايات المنقولة (جمعاً وتحرياً وتثبتاً ونقداً لمعرفة الصحيح)، ثم العناية بالفقه والفهم والتفسير. وكل العلوم الإسلامية إنما تدور حول كتاب الله عز وجل، فهو الأصل والحق المبين.

\subsection*{جهود أئمة الإسلام في حفظ وتصنيف هذه العلوم}
من هنا كانت العناية؛ فكل عالم من علماء الإسلام اعتنى بأن يخدم علوم الشريعة من باب من هذين البابين (الرواية أو الفقه). وقد صُنفت الكتب في جمع السنن والآثار، وكل إمام أراد أن يجمع هذه النصوص على طريقة تكمل هذا البناء العظيم. فنجد تنوعاً في التصنيفات:
\begin{itemize}
    \item \textbf{في باب الجمع العام:} الموطآت، المسانيد، الجوامع، المصنفات، السنن، المستخرجات، والمستدركات.
    \item \textbf{في باب النقد والرجال:} كتب الجرح والتعديل، كتب الرجال، كتب العلل، وكتب السؤالات.
    \item \textbf{في الأبواب المخصوصة:} مثل كتاب «الإيمان» لأبي عبيد، أو «خلق أفعال العباد» للبخاري في باب القدر.
    \item \textbf{في التفسير المسند:} مثل تفسير ابن أبي حاتم، وتفسير ابن جرير الطبري.
\end{itemize}

كل إمام من هؤلاء الأئمة أكمل لبنة في هذا البناء. ومن هذه الحلقات المهمة والمفصلية: \textbf{تفسير ابن جرير الطبري} رحمه الله. وأنا لا أحب أن ألقنك قيمة هذا الكتاب تلقيناً، وإنما ستعرفها وتلمسها أنت بنفسك؛ وهذا من أخص مقاصد هذه الدورة: ألا نكثر من المقدمات النظرية التي يمكن أن يصل إليها الطالب وحده، بل الغرض الأسمى أن تكتشف أنت قيمة الكتاب ومنهجه بنفسك أثناء قراءتنا وتدارسنا له.

\section*{الخطوط الثلاثة الرئيسية في طلب العلم}
طالب العلم تتدرج به الحلقات والأبواب. وأحدثكم هنا عن تجربتي؛ ففي طلب العلم عندي خطوط ثلاثة رئيسية:

\begin{enumerate}
    \item \textbf{الخط الأول (المصادر الأصلية التأسيسية):} وهو الأساس، ويشمل المصادر الأصول في علوم الشريعة. وخلاصة هذه المصادر تبدأ من كتاب «الموطأ» للإمام مالك إلى كتاب «المجتبى» للنسائي. هذه الطبقة هي أهم طبقة في العلوم؛ وتشمل: «الأم» للشافعي، و«الآثار» لمحمد بن الحسن الشيباني، ومصنف ابن أبي شيبة، ومصنف عبد الرزاق، وصحيح البخاري، وصحيح مسلم، وسنن الترمذي، والنسائي، وأبي داود، وابن ماجه، والدارمي، ومسند أبي داود الطيالسي. وكذلك في التفسير: «تفسير ابن جرير الطبري» رحمه الله، وقبله «تفسير ابن أبي حاتم».
    
    هذه الكتب هي أهم ما ينبغي أن يعتني به طالب العلم قراءةً ودراسةً وشرحاً واستخراجاً للفوائد. هي القاعدة والمتكأ، وكل العلوم أُخذت من هذه المصنفات وبُنيت عليها في كل أبواب الدين. فبقدر ما تعتني بهذا الخط قراءةً وضبطاً وفهماً، تنبل في هذا العلم إن شاء الله.

    \item \textbf{الخط الثاني (الأئمة المحققون من المتأخرين):} وهم الذين اعتنوا بجمع ما جاء في كل باب من أبواب العلم من السنن والآثار، مع شرح وفقه لهذا الجمع، وتصنيفه، واستخراج الفوائد منه. وأقصد بالمتأخرين هنا (بداية من القرن الرابع الهجري فصاعداً)؛ فإلى سنة 300 هـ تقريباً دونت أهم كتب الإسلام، وكان حقاً على من جاء بعدهم أن يجمع هذه العلوم ويضبطها ويبني عليها.
    
    من أمثلة هؤلاء الأئمة المحررين: ابن عبد البر (الذي اعتنى بموطأ مالك في كتابه «التمهيد»)، وابن الصلاح، وابن تيمية، وابن القيم، وابن حجر، وابن كثير، والشاطبي، والقرافي، وابن رجب الحنبلي. هؤلاء قصد كل واحد منهم علماً أو أكثر، فاعتنى بجمعه وضبطه وتحريره والتصنيف فيه.

    \item \textbf{الخط الثالث (الباحثون والدارسون المتميزون من المعاصرين):} وهو خط مهم جداً رغم كونه أقل مما سبق من ناحية أولوية العناية. وهؤلاء المتخصصون المعاصرون اعتنوا بأبواب العلم، واستقرأوا فيها، وأخرجوا أبحاثاً ودراسات تخدم العلم وتكون مدخلاً لطالب العلم لفهم كتب الأئمة.
    
    من أمثلة هؤلاء: الشيخ الألباني، والشيخ ابن عثيمين، والشيخ عبد الله السعد، والشيخ محمد عمرو عبد اللطيف، والشيخ طارق عوض الله. وفي علوم القرآن: د. مساعد الطيار، وفي علم العقيدة والفقه: د. يوسف الغفيص، وغيرهم. وطالب العلم يحتاج إلى دروسهم ومؤلفاتهم لا لتكون هي الأساس الذي يقف عنده، بل لتكون مجرد مدخل يتصور به العلوم ليلج منها إلى كتب الأئمة المتقدمين.
\end{enumerate}

ولا أريد أن أطيل في هذه المقدمات، فقد سبقت لنا دورات تمهيدية تأسيسية لطالب العلم (مثل: دورة بناء طالب العلم، رفع كفاءة طالب العلم، تهيئة المسلم للفقه في الدين، معوقات النبوغ لطالب العلم، نصائح وتوجيهات للدارسين عن بُعد، ودورة دراسة الكتب التأسيسية). وأنصح المبتدئ قبل الدخول في دراسة الطبري أن يستمع لدورة «نصائح وتوجيهات للدارسين عن بُعد» ودورة «دراسة الكتب التأسيسية».

\separator

\section*{السياق الدراسي لدراسة تفسير الطبري}
حتى تتصور موقع دراستنا لتفسير الطبري من خارطة دراستك للعلم؛ فنحن ندرس المصنفات التأسيسية في علوم الشريعة، ونريد أن يكون لنا عناية بأهم كتب الإسلام تحت كل باب. وقد درسنا ولله الحمد عدة كتب تأسيسية ككتاب «الرسالة» للشافعي، ومقدمة الإمام مسلم، وصحيح الإمام البخاري، ورسائل ابن تيمية وغيرها.

إذن، دراستنا في هذه الدورة ليست مجرد دراسة لتفسير القرآن أو علوم القرآن فحسب؛ بل نحن ندرس أحد أهم مصنفات الإسلام، وهو \textbf{تفسير ابن جرير الطبري} رحمه الله. هذا المصنف هو كنز من كنوز العلم من حيث الجمع، والتحرير، والنقد، والبيان، والاستدلال. ولا شك أن من أهم ما فيه هو تفسير القرآن، ومن أهم ما فيه أن تعرف كيف فُسر القرآن، ولماذا يختلف المفسرون، وما أصناف اختلافهم، وما هو منهج الطبري في التفسير. كل هذا محل اهتمام بلا شك، لكنه ليس الهدف الوحيد.

هذا الكتاب هو خلاصة رحلة الإمام الطبري رحمه الله، وهو زبدة تحصيله في علوم الإسلام، في علوم القرآن والحديث والآثار والفقه والإيمان وغير ذلك. فلم تكن دراستي له يوماً ضمن دراستي لتفسير القرآن فقط، بل ضمن دراستي لتراث \textbf{الأئمة المحققين} الذين هداني الله للعناية بسيرهم وتكوينهم ومناهجهم.

\section*{الفرق بين العالم والإمام المحقق}
كثيراً ما أذكر مصطلح \textbf{«الأئمة المحققون»}. التحقيق معناه: إعطاء الأمر حقه. فأهل العلم درجات، وليس كل من تكلم في العلم يكون محققاً؛ فقد يكون بعضهم جامعاً فقط، أو مجرد ناقل. لكن المحققين لهم صفات تميزهم؛ فهم الأئمة الذين حققوا ما عُنوا به من أبواب العلم ومسائله، ولهم آثار ظاهرة في تلقي العلوم وجمعها وضبطها وتحريرها والتأثير فيها تدريساً وتصنيفاً ونقداً وتأصيلاً. وصارت مصنفاتهم إماماً لمن بعدهم، فكل من يدخل في العلم لابد أن يلج من هذه المصنفات.

\subsection*{هل أعتني بآثار كل من صنف في العلم؟}
لا طبعاً. أهل العلم متفاوتون، ومسمى \textbf{«العالم»} يندرج تحته كل من تعاطى العلم وجدَّ فيه وحصل قدراً منه، وكان على علم بأبواب العلم ومسائله ودلائله ومواضع الإجماع والخلاف. 

أما \textbf{«الإمام المحقق»} فهو الذي جمع مع ما تقدم \textbf{التأثير في العلم}، أي له أثر بارز بحيث لا يُذكر العلم إلا ويُذكر هو. وهذا لا يكون أبداً بمجرد أن يكون جامعاً ناقلاً؛ بل لابد أن يكون فيه تحرٍّ وضبط، وأن يتعلم كل العلوم ليخدم العلم الذي يطلبه. 

ومن علامات هؤلاء الأئمة البارزة: \textbf{أن الإمام المحقق إذا فاتك يَصعُب أن تعوّضه.} 
\begin{itemize}
    \item من أراد فقه الشريعة وفاته «الشافعي» لا يجبره غيره.
    \item ومن أراد التفسير وفاته «الطبري» فاته أمر لا ينجبر.
    \item ومن فاتته الكتب التأسيسية التي بُني عليها العلم (كموطأ مالك، ومسند أحمد، ومصنف عبد الرزاق، ومصنف ابن أبي شيبة)، كيف يعوضها؟
\end{itemize}
هذه الكتب لم تؤلف لتقرأ وتُنسى، بل أُلفت لتقرأ وتُراجع وتُذاكر ويُشرح ويُبنى عليها. هي كتب ينبغي أن تكون ورداً يومياً تعيش عليه، ولا يصح أن نضيع أعمارنا في الحواشي والمختصرات ونطوف حول الكتب الأصلية كأنها أضرحة دون أن ندخل فيها.

\separator

\section{كيف نعرف أئمة كل باب؟}
كثيراً ما يسأل الطالب: كيف أعرف من هم أئمة الحديث؟ أو أئمة الجرح والتعديل؟ أو أئمة الفقه؟ وما هي أهم مصنفاتهم؟ 

هذا الأمر تعرفه بأحد طريقين:
\begin{enumerate}
    \item \textbf{التلقي والسؤال:} بأن تسأل من سبقك من المتخصصين في العلم ليدلك على أهم العلماء وأهم الكتب في بابه.
    \item \textbf{الدخول في العلم والاستقراء:} وهو الأهم؛ فبدخولك في العلم وممارسته ستكتشف هؤلاء الأئمة بنفسك. فمثلاً:
    \begin{itemize}
        \item لو تتبعت كتب الجوامع والسنن، ستعرف من هم أئمة الفقه من لدن الصحابة إلى نهاية القرن الثالث الهجري.
        \item ولو اعتنيت بكتب التفسير، ستجد تكرار النقل عن أئمة ومفسرين بعينهم (مثل: \rawi{ابن عباس}، \rawi{مجاهد}، \rawi{سعيد بن جبير}، \rawi{طاووس}، وغيرهم).
        \item وفي علوم الحديث والنقد، ستجد الاستقراء يقودك لأسماء بعينها تتكرر دائماً (مثل: \rawi{شعبة}، \rawi{يحيى القطان}، \rawi{عبد الرحمن بن مهدي}، \rawi{أحمد بن حنبل}، \rawi{يحيى بن معين}، و\rawi{علي بن المديني}).
    \end{itemize}
\end{enumerate}

والمشكلة أن كثيراً من الطلاب يريد أن يعرف كل شيء عن العلم وخطته قبل أن يدخل فيه! وهذا لن تعرفه حق المعرفة إلا بدخولك في العلم.

\section{أهمية دراسة سير الأئمة}
من الأمور التي أؤكد عليها كثيراً: دراسة سير الأئمة. هذا مكوّن أساسي لا يعوضه غيره أبداً في بناء طالب العلم. والترتيب المنهجي في ذلك: دراسة سير الأنبياء عليهم السلام، ثم سير المصلحين المذكورين في القرآن، ثم سيرة الصحابة رضي الله عنهم، ثم سيرة علماء الأمة.

\subsection{ما الذي نعتني به عند دراسة سير الأئمة؟}
حين تدرس ترجمة إمام من الأئمة، يجب أن تركز على استخراج الفوائد التالية:
\begin{enumerate}
    \item \textbf{النشأة والتكوين:} كيف نشأوا؟ ما هي العلوم التي درسوها؟ وكيف رحلوا في طلب العلم وتلقوه وصنفوه؟
    \item \textbf{معرفة قدر النفس:} بقراءتك لسير هؤلاء ستعرف قدر نفسك، فلا يغرك مدح المادحين حين ترى أين أنت من جَلَدِ هؤلاء وصبرهم واحتسابهم.
    \item \textbf{التمييز بين الأصول والفضول:} لتعرف ما هي كبار المسائل التي اشتغلوا بها؛ فكثير من الطلاب يعيش سنوات طويلة على لطائف خفيفة أو فضول مسائل ويعمى عن الأصول. بدراستك لسيرهم ستعرف أنهم أفنوا أعمارهم في كبار المسائل.
    \item \textbf{فهم سياق عصرهم:} العلم بواقعهم وأهم المشكلات والقضايا التي عالجوها. فكل إمام كان عالماً بعصره وقدم جواباً لما يحتاجه الناس؛ فالشافعي صنف «الرسالة» لحاجة عصره، والبخاري صنف «الجامع الصحيح» لحاجة عصره، والدارقطني صنف «العلل»، والطبري كذلك. فمعرفة عصر الإمام تجعلك تفهم لماذا وكيف صنف كتابه.
\end{enumerate}

\separator

\section{التأهب لقراءة تفسير الطبري والمطولات}
هذا الصنف من الكتب الذي نتكلم عنه، مثل «تفسير الطبري» أو «منهاج السنة» لابن تيمية، أو «الأم» للشافعي، أو «صحيح البخاري»، يحتاج إلى تأهب واستعداد. لا يصح للطالب أن يدخل في هذه الكتب المطولة دون مقدمات.

\subsection{خطورة الدخول في المطولات دون مقدمات}
في بداية دراستي لم يرشدني أحد إلى هذا، فدخلت في بعض هذه الكتب فوقعت في نوعين من الخطأ:
\begin{enumerate}
    \item \textbf{عدم فهم الكلام:} أقرأ في كتاب كـ «منهاج السنة» أو «درء تعارض العقل والنقل» أو «صحيح البخاري» أو «تفسير الطبري» فلا أفهم ماذا يقول العالم؛ لأنني لم أتدرب على طريقته وألفاظه ومصطلحاته.
    \item \textbf{فوات كنوز من الفوائد:} تمر عليّ فوائد عظيمة في أبواب التفسير، والسنن، والفرق، ولا أقف عليها أو لا أظن أنها فائدة أصلاً لغياب التصور المسبق للعلم.
\end{enumerate}

\subsection{أهمية المقدمات التأهيلية}
من هنا وجدت أن طالب العلم إذا أراد أن يدخل على كتاب من هذه الكتب التأسيسية فلابد أن يكون عنده تأهب. وهذه المقدمات من شأنها أمران:
\begin{itemize}
    \item \textbf{أن تعينك على فهم الكتاب:} الكتاب كالبناء الكبير يحتاج إلى أسس.
    \item \textbf{أن تكون كشافاً للفوائد:} تنبهك إلى أن هذه مسألة، وهذه فائدة، وهذه قاعدة، وهذا إشكال وجوابه. هي كالكشاف الذي ينير لك الطريق لتستخرج الفوائد.
\end{itemize}

\subsection{الكتب المقترحة للتأهب قبل قراءة الطبري}
تفسير الطبري بحر في علوم القرآن والتفسير والحديث والفرق والفقه واللغة. وحتى نفهم هذا الكتاب ونستخرج أقصى ما يمكن من الفوائد، نحتاج إلى أرضية تناسب موضوعه. وهذه قائمة مقترحة لتكون مدخلاً لك لتتصور علوم القرآن:

\begin{itemize}
    \item \textbf{في أصول التفسير ومقدماته:} «مقدمة في أصول التفسير» لابن تيمية (بشرح د. مساعد الطيار)، «التحرير في أصول التفسير» لد. مساعد الطيار، و«مختصر قواعد الترجيح عند المفسرين» لحسين الحربي.
    \item \textbf{في علوم القرآن:} «المحرر في علوم القرآن» لد. مساعد الطيار، و«المقدمات الأساسية في علوم القرآن» لعبد الله الجديع.
    \item \textbf{في غريب القرآن (معاني الكلمات):} «السراج في بيان غريب القرآن» لمحمد الخضيري، أو «وجه النهار الكاشف عن معاني كلام الواحد القهار» لعبد العزيز الحربي.
    \item \textbf{في التفاسير المختصرة:} «المختصر في التفسير» (مركز تفسير)، أو «المعين في تدبر الكتاب المبين» لمجد مكي.
    \item \textbf{في مناهج المفسرين:} «اتجاهات التفسير في القرن الرابع عشر» لد. فهد الرومي، و«مختصر مناهج المفسرين» لمحمد أبو زيد.
    \item \textbf{أبحاث ودراسات نوعية مهمة:}
    \begin{itemize}
        \item «الأثر العقدي في تعدد التوجيه الإعرابي لآيات القرآن» لعبد الله السيف.
        \item «التفسير اللغوي للقرآن» لد. مساعد الطيار.
        \item «مفهوم التفسير والتأويل والاستنباط» لد. مساعد الطيار.
        \item «اختلاف السلف في التفسير» لمحمد صالح سليمان.
    \end{itemize}
\end{itemize}

وليس المقصود هنا ألا تبدأ في دراسة الطبري حتى تنجز وتنهي كل هذه الكتب! بل المراد أن هذه الكتب هي مدخل تتصور به العلم لتعرف أبوابه ومسائله، ولتكون لك أرضية تبصرك بما تطلبه من الفوائد في تفسير الطبري. ولتيسير الأمر في دورتنا هذه، اكتفينا بثلاثة كتب كمدخل: (مقدمة التفسير لابن تيمية، التحرير في أصول التفسير، والمحرر في علوم القرآن).

\separator

\section*{الخطوات العملية لدراسة تفسير الطبري}
بعد أن تحدثنا عن أهمية دراسة تفسير الطبري، وعن الخطوط الرئيسية في طلب العلم، وعن أهمية دراسة سير الأئمة، وعن التأهب لقراءة المطولات، نأتي الآن إلى الخطوات العملية لدراسة تفسير الطبري. هذه الخطوات هي التي ستقودك إلى فهم هذا الكتاب واستيعابه واستخراج الفوائد منه.

\subsection*{خطوة 1: التحضير المسبق}
قبل أن تبدأ في دراسة تفسير الطبري، يجب أن تقوم بتحضير مسبق. هذا التحضير يشمل:
\begin{itemize}
    \item قراءة المقدمات التأهيلية التي ذكرناها سابقاً.
    \item الاطلاع على سير الإمام الطبري لتتصور شخصيته وعصره ومناهجه.
    \item تجهيز أدوات الدراسة مثل الدفتر والقلم، أو البرامج الإلكترونية إذا كنت تدرس عن بُعد.
\end{itemize}

\subsection*{خطوة 2: المشاركة الفعالة في الدروس}
عندما تبدأ في حضور دروس تفسير الطبري، يجب أن تكون مشاركاً فعالاً. هذا يعني:
\begin{itemize}
    \item الاستماع بتركيز والانتباه لما يقال.
    \item طرح الأسئلة والاستفسارات عند الحاجة.
    \item المشاركة في النقاشات والمناقشات الجماعية.
\end{itemize}

\subsection*{خطوة 3: المراجعة واستخراج الفوائد}
بعد كل درس، يجب أن تقوم بمراجعة ما تم دراسته واستخراج الفوائد. هذا يشمل:
\begin{itemize}
    \item كتابة الفوائد التي استخلصتها من الدرس وتصنيفها تحت بابها اللائق.
    \item جمع النظائر والشواهد والأمثلة التي تم ذكرها في الدرس.
    \item محاولة تطبيق ما تعلمته في فهم آيات أخرى أو في مسائل أخرى.
\end{itemize}

\separator

\section{عناوين الفوائد المستخرجة من تفسير الطبري}
من خلال قراءتي لكتاب الطبري، حصرت لكم عناوين مهمة جداً هي التي سنطلبها ونكلف الطلاب بجمعها. وكلما نقف على فائدة منها سنقول: هذه فائدة في كذا، بحيث يأخذ كل واحد منكم ملفاً من هذه الفوائد ويجمع فيه الفوائد المتعلقة به.

أهمية ذلك أن هذه الكتب غالباً ما تُقرأ في العمر مرة واحدة؛ لأن العمر لا يتسع. لذلك، إذا دخلت على كتاب فضع في ذهنك أنك لن تقرأه مرة أخرى؛ فهذا أدعى إلى أن تفهمه وتستخرج منه كل فائدة ولا تزهد فيها.

أما عناوين الفوائد التي سنجمعها فهي:
\begin{itemize}
    \item \textbf{منهج الطبري في التفسير:} كيف يفسر، ومصادره في التفسير واللغة، والقواعد التي يذكرها.
    \item \textbf{منهجه في عرض الأقوال والخلاف:} كيف يعرض الخلاف، وبمن يبدأ، وأسباب الخلاف في التفسير وأنواعه (كالقراءات، أسباب النزول، الإعراب، عودة الضمير، ومعاني الألفاظ).
    \item \textbf{موقفه من الإسرائيليات والقراءات:} لماذا يوردها؟ وهل يقبلها أم يردها؟ وعمن يرويها؟ وما ضابطه في قبول أو رد أو ترجيح القراءات؟
    \item \textbf{منهجه في الاستدلال وطرح الإشكالات:} كيف يستدل للقول الذي يرجحه؟ وكيف يطرح الإشكال ويجيب عنه؟ وما هي قرائن الترجيح عنده؟
    \item \textbf{الاستنباطات:} منهجه في استخراج المعاني الدقيقة من الآيات.
    \item \textbf{أقواله في الاعتقاد:} رأيه في أبواب الإيمان، أسماء الله وصفاته وأفعاله، النبوات، القدر، ونقده للفرق والمقالات (كالمرجئة، الخوارج، والمعتزلة).
    \item \textbf{المفسرون من السلف وأسانيده إليهم:} من هم المعتبرون عنده من الصحابة والتابعين؟ وما هي أهم أسانيده إليهم؟ وكيف يستدرك عليهم ويناقش أقوالهم؟
    \item \textbf{الأمثلة والشواهد التطبيقية:}
    \begin{itemize}
        \item أمثلة لتفسير القرآن بالقرآن، وبالحديث، وبأقوال الصحابة، وبروايات بني إسرائيل.
        \item أمثلة للنسخ.
        \item الشواهد الشعرية (من هم الشعراء؟ ما هي الطبقة التي يستشهد بها؟ وعلى ماذا يستشهد؟).
    \end{itemize}
    \item \textbf{الفوائد اللغوية وعلوم القرآن والحديث:} استدراكه على اللغويين (كأبي عبيدة)، وفوائد في المكي والمدني وجمع القرآن، وفوائد في نقد الأسانيد وعلل الروايات، وغريب القرآن (الألفاظ التي فسرها الطبري).
    \item \textbf{إبداع الطبري:} ماذا أضاف الطبري لمن سبقه في التفسير؟ وما هي أوجه تميزه؟
\end{itemize}

\separator

\section{خطة البدء في قراءة الكتاب وخاتمة اللقاء}
بعد هذه المقدمة، خطتنا كالتالي: سنبدأ أولاً بقراءة مقدمة تفسير الطبري، وسنعتمد على شرح الدكتور مساعد الطيار لها. بعد الانتهاء منها، سندخل في قراءة التفسير كلمة كلمة، حتى الأسانيد سنقرأها. 

وسيكون عملي في الكتاب معكم يتركز في أمرين:
\begin{enumerate}
    \item \textbf{الفهم:} فك العبارات التي أراها صعبة حتى نفهم مراد الإمام الطبري.
    \item \textbf{الاستخراج:} أن أدلك على الفوائد وكيف توظفها وتقيدها تحت أبوابها التي ذكرناها.
\end{enumerate}

أرجو أن يكون هذا الكتاب مفتاحاً لك من الخير. فهذا الكتاب حينما قرأته قبل سنوات شعرت أنني وُلدت من جديد؛ ليس في التفسير فقط، بل في علوم الشريعة بأكملها وفي طريقة تناول المسائل. 

أرجو أن تعتنوا وتستعينوا بالله، وتكثروا من الدعاء أن يعيننا الله ويبارك في أوقاتنا ويتم لنا هذا الخير. بارك الله فيكم وجزاكم الله خيراً، والسلام عليكم ورحمة الله وبركاته.